\begin{enabstract}
With the rapid development of low Earth orbit (LEO) satellite communication systems and non-terrestrial networks (NTN), direct-to-smartphone satellite access is evolving from emergency communication toward broader direct-to-terminal connectivity. While this trend improves wide-area access capability and communication flexibility, it also brings new challenges to radio security supervision. In complex urban environments, non-cooperative satellite direct-connect terminals may exploit building blockage and complicated propagation conditions to conduct concealed access or unauthorized radiation, while abnormal clues usually provide only coarse suspicious areas and cannot directly determine the number and spatial positions of targets. Compared with conventional fixed monitoring stations, low-cost unmanned aerial vehicle (UAV) platforms provide higher mobility, more flexible deployment, and adjustable observation perspectives, and therefore offer a new technical approach for non-cooperative terminal localization and investigation in dense urban areas. However, under the combined influence of urban multipath, non-line-of-sight propagation, local blockage, and the limitations of platform payload and onboard computing capability, traditional geometric localization methods and static fusion strategies often fail to balance localization accuracy, environmental adaptability, and engineering practicability. To address these issues, this thesis studies non-cooperative localization of satellite direct-connect terminals from two aspects: single-target continuous localization and multi-target robust cooperative localization.

For single-target continuous localization, the reliability of multi-source observations varies with the propagation environment and motion state, making fixed-weight fusion difficult to adapt stably to complex scenarios. To address this issue, a physics-aware deep reinforcement learning based adaptive fusion localization method is proposed. The method establishes a weighted fusion model for localization results from received signal strength indicator (RSSI), angle of arrival (AOA), and pedestrian dead reckoning (PDR), formulates the dynamic weight adjustment process as a sequential decision problem, and introduces state features including root mean square (RMS) delay spread, RSSI variance, geometric consistency, and historical fusion weights. A Physics-Aware Deep Q-Network (PA-DQN) agent is further designed to adaptively adjust fusion weights through incremental actions and a reward function considering both localization accuracy and weight smoothness. Experimental results show that, on the overall test set, the proposed method achieves a root mean square error (RMSE) of 2.15 m and a 95th-percentile error (P95) of 3.87 m, and outperforms RL-IFA and EKF in overall localization accuracy, tail-error control, low-SNR robustness, and trajectory smoothness.

For multi-target concurrent localization, uncertainty in the number of targets, sparse UAV spatial scanning, and multipath-induced false peaks make traditional peak-based decision methods unreliable. To address this challenge, a robust cooperative localization method based on spatial signal representation is further proposed. The method constructs a complex urban propagation field with OpenStreetMap (OSM), a MATLAB toolchain, and ray-tracing modeling, and transforms discrete spatial observations into structured signal representations through RSSI grid mapping and Radio Map reconstruction. For random node failures in multi-UAV cooperative networks, a teacher-student knowledge distillation framework is designed, where the teacher model under complete observations supervises the student model under missing observations, enabling the student network to approximate the decision capability of the complete network. Experimental results demonstrate that the proposed method effectively alleviates performance degradation caused by random node failures, maintains relatively stable localization performance under medium and high failure ratios, and improves robust cooperative localization capability in complex environments.

In summary, this thesis establishes a localization and identification framework for non-cooperative satellite direct-connect terminals based on low-cost UAV platforms. The results demonstrate that integrating physical priors, deep reinforcement learning, and knowledge distillation improves localization accuracy and robustness in complex environments, providing useful references for radio security supervision in complex scenarios.

\enkeywords{UAV; Satellite Communication; Fusion Localization; Radio Map; Knowledge Distillation}
\end{enabstract}
